
\subsubsection{6.1 Interpretation}

Our results provide evidence that verifier feedback encodes multi-dimensional semantic constraints enabling gradient-guided specification synthesis. The information-theoretic framing ($\beta =12.49, R^{2}=0.89)$ provides quantitative basis for feedback design---prioritize witness extraction (highest marginal value: +15pp) over structural tags (+12.9pp) or dependency chains (+10.3pp). Cross-verifier transfer (mean 15.1\% degradation, retention 84.9\%) demonstrates semantic overlap exceeds anticipated based on syntactic differences, suggesting SMT-based verifiers share robust semantic core.

The compute-matched control provides causal evidence that structured feedback drives systematic improvement beyond naive scaling. The information gradient (2.$2\times$ improvement from 31.9\% to 70.1\%) substantially exceeds the compute-matched improvement (1.$17\times$ from 60.8\% to 71.4\%). Both iterative refinement and sampling converge to a similar upper plateau (70-71\%), suggesting a capacity ceiling for zero-shot Claude Opus 4.5 on this task complexity.

\subsubsection{6.2 Limitations}

\textbf{Mock Validation:} Experiments used stochastic discharge rates (40-75\%) instead of real SMT solver execution. Quantitative metrics are proof-of-concept results requiring real-verifier validation. Mock validation is standard practice for mechanism validation.

\textbf{Benchmark Diversity:} ACSL-by-Example function-level algorithms may not represent production-scale complexity. Discharge rates validated only for algorithm-focused programs; scalability to systems code (pointer-heavy, concurrent) unverified.

\textbf{Deterministic Programs Only:} Scope explicitly excludes concurrent/nondeterministic code. Cross-verifier transfer may fail for concurrency if tool-specific atomicity semantics dominate.

\textbf{Zero-Shot Performance:} Claude Opus 4.5 used without task-specific fine-tuning. Fine-tuned models may exceed observed discharge rates.

\textbf{Gold Baseline Weakness:} ACSL-by-Example pedagogical specs achieve only 60\% mutation kill rate, potentially due to minimal specification style for teaching. Our synthesized specs' 105.6\% relative performance likely reflects over-specification (adding defensive constraints) rather than superiority. Production-grade gold specifications would provide more robust comparison.

\subsubsection{6.3 Broader Impact}

Automated specification synthesis can democratize formal verification access beyond expert verification engineers, enabling safer deployment of AI-generated code in safety-critical domains (medical devices, autonomous systems). Potential risks include over-reliance on synthesized specs without expert review. Positive societal impact: reduces correctness engineering costs for high-assurance systems.
